% U5 WS4 -- pre-equilibrium, multistep profiles, catalysis + FRQ. Blocks 6-7.
\documentclass{apchem-worksheet}

\apcunit{5}
\apcunittitle{Kinetics}
\apcdoctitle{Worksheet 4 \textbullet\ Pre-Equilibrium, Catalysis \& FRQ}
\apcsubtitle{CED 5.9--5.11 \textbullet\ a final rate law never contains an intermediate}

\begin{document}
\wsheader[Fast first step: write the slow step's rate law, then use the
equilibrium of step 1 to eliminate the intermediate \textbullet\ a catalyst
lowers $E_a$ for \emph{both} directions and changes neither $\Delta H$
nor $K$.]

\problem[]{\ced{5.9} Overall $\ce{2NO + O2 -> 2NO2}$, experimental
rate $= k[\ce{NO}]^2[\ce{O2}]$. Mechanism: step 1 (fast equilibrium)
$\ce{2NO <=> N2O2}$; step 2 (slow) $\ce{N2O2 + O2 -> 2NO2}$.}
\begin{parts}
  \item Write the rate law for the slow step as it stands.
        \answerblank[5cm]{rate $= k_2[\ce{N2O2}][\ce{O2}]$}
  \item Why is that not yet an acceptable rate law?
        \answerlines[2]{\ce{N2O2} is an intermediate. Its concentration is
        not something you can measure or control, so it cannot appear in a
        final rate law.}
  \item Use step 1's equilibrium to express $[\ce{N2O2}]$.
        \answerblank[5cm]{$[\ce{N2O2}] = K_1[\ce{NO}]^2$}
  \item Substitute and give the predicted rate law.
        \workspace[2cm]
        \keyonly{\begin{apcanswer}rate $= k_2K_1[\ce{NO}]^2[\ce{O2}] =
        \mathbf{k[\ce{NO}]^2[\ce{O2}]}$, matching experiment, with the
        observed $k = k_2K_1$\end{apcanswer}}
  \item What does it mean that the observed $k$ equals $k_2K_1$?
        \answerlines[3]{The measured rate constant is a composite of an
        elementary rate constant and an equilibrium constant, not a property
        of any single step. That is one reason the same rate law can arise
        from more than one mechanism.}
\end{parts}

\problem[]{\ced{5.9} Derive the rate law for: $\ce{A + B <=> C}$ (fast
equilibrium), then $\ce{C + B -> D}$ (slow). \workspace[3cm]
\keyonly{\begin{apcanswer}Slow step: rate $= k_2[\ce{C}][\ce{B}]$. From
step 1, $K_1 = [\ce{C}]/([\ce{A}][\ce{B}])$, so
$[\ce{C}] = K_1[\ce{A}][\ce{B}]$. Substituting gives
rate $= k_2K_1[\ce{A}][\ce{B}]^2 =
\mathbf{k[\ce{A}][\ce{B}]^2}$ --- note B ends up \emph{second}
order.\end{apcanswer}}}

\problem[]{\ced{5.10} Multistep energy profiles.}
\begin{parts}
  \item The number of humps equals the number of:
        \answerblank[3.4cm]{elementary steps}
  \item A valley between two humps is:
        \answerblank[3.2cm]{an intermediate}
  \item The rate-determining step is the one with:
        \answerblank[6.2cm]{the highest activation barrier}
  \item A profile shows a first barrier of \SI{95}{\kilo\joule} and a
        second of \SI{60}{\kilo\joule} measured from the intermediate.
        Which step is rate-determining?
        \answerblank[2.6cm]{the first}
  \item If instead the second barrier were the larger, what would change
        about the rate law? \answerlines[3]{The rate law would come from
        the second step. Since that step consumes an intermediate, the
        intermediate would have to be eliminated using the equilibrium of
        the first step, and the resulting rate law would generally look
        different.}
\end{parts}

\problem[]{\ced{5.11} Catalysis.}
\begin{parts}
  \item A catalyst increases the rate by providing:
        \answerblank[7.5cm]{an alternative pathway with lower $E_a$}
  \item List three quantities a catalyst does \emph{not} change.
        \answerblank[6cm]{$\Delta H$, $K$, $\Delta G^\circ$}
  \item Explain why a catalyst cannot shift an equilibrium.
        \answerlines[3]{It lowers the activation energy of the forward and
        reverse reactions by the same amount, so both rates increase by the
        same factor. The ratio that defines the equilibrium position is
        unchanged --- the system simply arrives there sooner.}
  \item Distinguish homogeneous from heterogeneous catalysis.
        \answerblank[8.7cm]{same phase as reactants vs different phase}
\end{parts}

\problem[]{\ced{5.11} For the mechanism $\ce{H2O2 + I- -> H2O + IO-}$, then
$\ce{H2O2 + IO- -> H2O + O2 + I-}$:}
\begin{parts}
  \item Overall equation:
        \answerblank[5cm]{$\ce{2H2O2 -> 2H2O + O2}$}
  \item Catalyst: \answerblank[2.4cm]{\ce{I-}}
  \item Intermediate: \answerblank[2.4cm]{\ce{IO-}}
  \item State the test that distinguishes them.
        \answerlines[2]{A catalyst is consumed in an early step and
        regenerated later, so it is present at the end. An intermediate is
        produced first and consumed later, so it is absent at the end and
        from the overall equation.}
\end{parts}

\problem[]{\textbf{FRQ (10 points).} \ced{5.2} \ced{5.3} \ced{5.8}
The gas-phase reaction $\ce{2NO(g) + O2(g) -> 2NO2(g)}$ was studied at
constant temperature.}
\begin{center}
\begin{tabular}{@{}cccc@{}}
\hline
\textbf{Exp} & $[\ce{NO}]_0$ (M) & $[\ce{O2}]_0$ (M) &
  \textbf{initial rate} (\si{\Molar\per\second}) \\
\hline
1 & 0.010 & 0.010 & $2.5\times10^{-5}$ \\
2 & 0.020 & 0.010 & $1.0\times10^{-4}$ \\
3 & 0.010 & 0.020 & $5.0\times10^{-5}$ \\
\hline
\end{tabular}
\end{center}
\begin{parts}
  \item Determine the order in \ce{NO} and in \ce{O2}, showing the
        comparison used for each. \answerlines[3]{Comparing 1 and 2, only
        $[\ce{NO}]$ changed: doubling it multiplied the rate by 4, so the
        order in \ce{NO} is 2. Comparing 1 and 3, doubling $[\ce{O2}]$
        doubled the rate, so the order in \ce{O2} is 1.}
  \item Write the rate law and calculate $k$ with units.
        \workspace[2.4cm]
        \keyonly{\begin{apcanswer}rate $= k[\ce{NO}]^2[\ce{O2}]$;
        $k = 2.5\times10^{-5}/[(0.010)^2(0.010)] =
        \mathbf{25}~\si{\per\Molar\squared\per\second}$\end{apcanswer}}
  \item Predict the initial rate when $[\ce{NO}] = 0.030$ and
        $[\ce{O2}] = 0.015$. \workspace[2.2cm]
        \keyonly{\begin{apcanswer}$(25)(0.030)^2(0.015) =
        (25)(9.0\times10^{-4})(0.015) =
        \mathbf{3.4\times10^{-4}}~\si{\Molar\per\second}$\end{apcanswer}}
  \item A student proposes that the reaction occurs in a single
        termolecular step. State whether this is consistent with the rate
        law, and whether it is likely. \answerlines[3]{It is consistent ---
        a single step $\ce{2NO + O2 -> 2NO2}$ would predict exactly
        rate $= k[\ce{NO}]^2[\ce{O2}]$. But it is unlikely, because it
        requires three particles to collide simultaneously with correct
        orientation, which is improbable.}
  \item Show that the mechanism $\ce{2NO <=> N2O2}$ (fast equilibrium)
        followed by $\ce{N2O2 + O2 -> 2NO2}$ (slow) predicts the same rate
        law. \workspace[2.8cm]
        \keyonly{\begin{apcanswer}Slow step: rate $= k_2[\ce{N2O2}][\ce{O2}]$.
        From the equilibrium, $[\ce{N2O2}] = K_1[\ce{NO}]^2$. Substituting:
        rate $= k_2K_1[\ce{NO}]^2[\ce{O2}]$, which has the observed
        form.\end{apcanswer}}
\end{parts}

\begin{apcnote}[Rubric --- score yourself, 10 points]
\textbf{(a) 3 pts} 1 order in \ce{NO} $=2$, 1 order in \ce{O2} $=1$, 1 for
citing the specific experiments compared --- stating orders with no
comparison earns 1 of 3 \textbullet\ \textbf{(b) 2 pts} 1 rate law, 1 value
of $k$ \emph{with units} \textbullet\ \textbf{(c) 2 pts} 1 correct
substitution, 1 answer $3.4\times10^{-4}$ \textbullet\ \textbf{(d) 1 pt}
consistent but improbable, with the termolecular reason \textbullet\
\textbf{(e) 2 pts} 1 writes the slow step's rate law, 1 eliminates the
intermediate using $K_1$ --- a final answer still containing \ce{N2O2}
earns 0 for the second point.
\end{apcnote}

\begin{spiralstrip}{Unit 5 \textbullet\ blocks 1--5}
  \item A final rate law must not contain:
        \answerblank[3.2cm]{an intermediate}
  \item Molecularity gives order only for:
        \answerblank[3.8cm]{an elementary step}
  \item $E_a(\text{fwd}) - E_a(\text{rev}) = $
        \answerblank[1.8cm]{$\Delta H$}
  \item Larger $E_a$ means temperature sensitivity is:
        \answerblank[2.4cm]{greater}
  \item Three requirements for a reactive collision:
        \answerblank[5.8cm]{contact, energy, orientation}
\end{spiralstrip}

\end{document}
